\section{整环}

记号约定:$A$是交换环,$B$是环.

\begin{definition}{整环}{}
	设$A$是非零环.若$A$没有零因子,则称$A$是整环.
\end{definition}

\begin{example}{}{}
	$\Z$是整环.
\end{example}

\begin{proposition}{域是整环}{}
	若$A$是域,则$A$是整环.
\end{proposition}

\begin{proposition}{}{}
	若$A$是整环,则$A$的每个子环都是整环.
\end{proposition}

\begin{corollary}{}{}
	若$A$是域,则$A$的每个子环都是整环.
\end{corollary}

\begin{proposition}{}{}
	若$A$是整环,则$A$的全分式环时域.
\end{proposition}

\begin{definition}{}{}
	设$A$是整环.我们称$A$的全分式环为$A$的分式域.
\end{definition}

\begin{proposition}{}{}
	设$B$是非零环,$A$是$B$的子环,$A\setminus\left\{0\right\}\subseteq B^{\times}$.
	\begin{enumerate}
		\item $A$是整环.
		\item 记$A^{\prime}$为$A$的分式域,则典范嵌入$A\hookrightarrow B$可唯一延拓为$f\in\Isom_{\mathsf{Ring}}\left(A^{\prime},B\right)$,其中$\im\left(f\right)$是$B$的子除环且
		\begin{align*}
			\im\left(f\right)=\left\{xy\in A\middle|x\in A,y\in A\setminus\left\{0\right\}\right\}.
		\end{align*}
	\end{enumerate}
\end{proposition}













